\documentclass{article}
\usepackage{amsmath, amssymb, amsthm}

\newtheorem{proposition}{Proposition}

\title{Extension of the Bellman Principle to Pareto Fronts in Multi-Objective Path Problems}
\author{}

\begin{document}
\maketitle

\section{Notation}

Let $G = (V, E)$ be a directed graph, where $V$ is the set of nodes and $E$ the set of edges. The state of the system is given by a node $x \in V$. The actions available at $x$ correspond to outgoing edges $e_{xj}$, where $j$ indexes the set of neighbors of $x$. Each edge $e_{xj}$ is assigned two non-negative costs: $d_{xj}$ (the first cost), and $p_{xj}$ (the second cost).

A \emph{path} from node $x_0$ to node $x_n$ is denoted by $\sigma_{x_0,x_n} = (e_{x_0 x_1}, e_{x_1 x_2}, \ldots, e_{x_{n-1} x_n})$.  
Composition of a path with an edge is written using the symbol $\circ$: for example,
\[
\sigma_{x_0, x_{n-1}} \circ e_{x_{n-1}, x_n} = \sigma_{x_0, x_n}.
\]

The \emph{cost vector} of a path $\sigma_{x_0, x_n}$ is defined by
\begin{equation}\label{eq:path-cost}
C(\sigma_{x_0, x_n}) = \left( \sum_{k=1}^n d_{x_{k-1} x_k},\; \sum_{k=1}^n p_{x_{k-1} x_k} \right).
\end{equation}

\section{Pareto Optimality and Value Set}

A cost vector $c = (d,p)$ is said to \emph{dominate} another vector $c' = (d',p')$ if $d \leq d'$ and $p \leq p'$ with at least one strict inequality.  
A path $\sigma_{x_0, x}$ is \emph{Pareto optimal} if there does not exist any other path $\sigma_{x_0, x}'$ from $x_0$ to $x$ with $C(\sigma_{x_0, x}')$ dominating $C(\sigma_{x_0, x})$.

The \textbf{value set} of node $x$ is defined as the set of cost vectors achieved by Pareto-optimal paths from $x_0$:
\begin{equation}\label{eq:pareto-set-def}
V(x) = \left\{ C(\sigma_{x_0, x}) : \sigma_{x_0, x} \text{ is Pareto optimal} \right\}.
\end{equation}

\section{Bellman Principle for Pareto Fronts}

We now state and prove the multi-objective extension of the Bellman principle for the value sets.

\begin{proposition}[Bellman Principle for Pareto Fronts]
Let $x_n$ be a node with incoming edges from the set of predecessors $x_{n-1}$ (i.e., $e_{x_{n-1}, x_n} \in E$). Then the value set at $x_n$ satisfies
\begin{equation}\label{eq:bellman-recursion}
V(x_n) = \operatorname{ParetoFront} \Bigg( 
  \bigcup_{x_{n-1}: e_{x_{n-1}, x_n} \in E} \left\{ \, c + (d_{x_{n-1}, x_n},\; p_{x_{n-1}, x_n}) \mid c \in V(x_{n-1}) \, \right\}
\Bigg),
\end{equation}
where $c + (d,p)$ denotes componentwise addition, and $\operatorname{ParetoFront}(\cdot)$ denotes the extraction of the set of non-dominated vectors.
\end{proposition}

\begin{proof}
Let us fix the terminal node $x_n$. By the definition of a path, any $\sigma_{x_0, x_n}$ must end with some predecessor $x_{n-1}$ and the edge $e_{x_{n-1}, x_n}$:
\begin{align*}
\sigma_{x_0, x_n} &= \sigma_{x_0, x_{n-1}} \circ e_{x_{n-1}, x_n}.
\end{align*}
By equation~\eqref{eq:path-cost}, the cost vector for such a path can be written as
\begin{equation}\label{eq:cost-recursion}
C(\sigma_{x_0, x_n}) = C(\sigma_{x_0, x_{n-1}}) + (d_{x_{n-1}, x_n},\; p_{x_{n-1}, x_n}).
\end{equation}
Therefore, the set of all cost vectors that can reach $x_n$ is the union over all predecessors $x_{n-1}$ and over all possible cost vectors $c$ from Pareto-optimal paths to $x_{n-1}$, with addition of the last edge's cost:
\begin{equation}\label{eq:all-candidates}
S(x_n) = \bigcup_{x_{n-1}: e_{x_{n-1}, x_n} \in E} \left\{ c + (d_{x_{n-1}, x_n},\; p_{x_{n-1}, x_n}) \mid c \in V(x_{n-1}) \right\}.
\end{equation}

The value set $V(x_n)$ must consist of those elements of $S(x_n)$ that are not dominated by any other element in $S(x_n)$, i.e., the Pareto front of $S(x_n)$. Indeed,  
\begin{itemize}
    \item If $c$ is Pareto-optimal at $x_{n-1}$, then $c + (d_{x_{n-1}, x_n}, p_{x_{n-1}, x_n})$ is a possible cost vector to $x_n$ via a path that is Pareto-optimal up to $x_{n-1}$.
    \item If a cost vector can be improved or matched in both components by some other extension from potentially another predecessor, then it is not Pareto-optimal at $x_n$ and is discarded.
\end{itemize}

Thus, to compute $V(x_n)$, it is sufficient and necessary to:
\begin{enumerate}
    \item Collect all cost vectors in $S(x_n)$ formed by extending all Pareto-optimal vectors from predecessors,
    \item extract the non-dominated subset, i.e., the Pareto front.
\end{enumerate}
This precisely matches equation~\eqref{eq:bellman-recursion}, establishing the recursive structure analogous to the classical Bellman principle.

\end{proof}

\textbf{Summary:}  
We have generalized the Bellman principle to the setting of multi-objective path problems, showing that the Pareto-optimal value set at each node is recursively constructed by extending the Pareto-optimal sets from all predecessors using the respective edge costs, followed by selection of the Pareto front. This provides a constructive and inductive characterization of optimal paths in multi-objective graph models.

\end{document}